\renewcommand{\arraystretch}{1.2}
\begin{tabular}{|l|r|r|r|r|>{\raggedright\arraybackslash}p{4.5cm}|}
\hline
\rowcolor{gray!20} \textbf{Arista} & \textbf{Oficial} & \textbf{Recta} & \textbf{Modelo} & \textbf{Razón} & \textbf{Causa} \\
\hline
\textbf{Chocontá -- Tunja} & 17.69 & 56.96 & 56.96 & 0.3106 & sector truncado en la fuente: el trazado queda a 43.70 km de un nodo \\
\hline
\textbf{Madrid -- Bogotá} & 9.12 & 19.91 & 19.91 & 0.4582 & entrada urbana a Bogotá: el trazado de INVÍAS empieza a 8.96 km del nodo \\
\hline
\textbf{Bogotá -- Tocancipá} & 22.12 & 41.19 & 41.19 & 0.5371 & entrada urbana a Bogotá: el trazado de INVÍAS empieza a 18.90 km del nodo \\
\hline
\textbf{Cajicá -- Zipaquirá} & 9.10 & 12.15 & 12.15 & 0.7492 & coordenada del nodo en el centro del municipio: el extremo del trazado queda a 2.14 km del nodo \\
\hline
\textbf{Ye de Granada -- Granada} & 2.43 & 3.04 & 3.04 & 0.7984 & coordenada del nodo derivada de la unión de tramos de INVÍAS (sin punto DANE) \\
\hline
\textbf{Girardot -- El Espinal} & 19.36 & 20.42 & 20.42 & 0.9479 & coordenada del nodo en el centro del municipio: el extremo del trazado queda a 3.07 km del nodo \\
\hline
\textbf{Bogotá -- Guasca} & 34.88 & 35.12 & 35.12 & 0.9931 & entrada urbana a Bogotá: el trazado de INVÍAS empieza a 11.17 km del nodo \\
\hline
\textbf{Tunja -- Duitama} & 47.73 & 47.78 & 47.78 & 0.9989 & coordenada del nodo en el centro del municipio: el extremo del trazado queda a 4.17 km del nodo \\
\hline
\end{tabular}
